% Bagian: computability
% Bab: recursive-functions
% Subbagian: non-pr-functions

\documentclass[../../../include/open-logic-section]{subfiles}

\begin{document}

\olfileid[id]{cmp}{rec}{npr}
\olsection{Fungsi yang Bukan Rekursif Primitif}

Fungsi-fungsi rekursif primitif tidak mencakup seluruh fungsi yang secara
intuitif dapat dihitung. Secara intuitif jelas bahwa kita dapat membuat daftar
semua fungsi rekursif primitif uner, $f_0$, $f_1$, $f_2$,~\dots, sedemikian
sehingga kita dapat menghitung secara efektif nilai $f_x$ pada masukan~$y$;
dengan kata lain, fungsi $g(x,y)$ yang didefinisikan oleh
\[
g(x,y) = f_x(y)
\]
dapat dihitung. Namun, fungsi
\begin{eqnarray*}
h(x) & = & g(x,x) + 1 \\
& = & f_x(x) +1
\end{eqnarray*}
juga dapat dihitung. Bagi setiap fungsi rekursif primitif $f_i$, nilai $h$ dan
$f_i$ berbeda pada~$i$. Jadi, $h$ dapat dihitung tetapi bukan rekursif
primitif; hal yang sama dapat dikatakan tentang~$g$. Ini merupakan versi
``efektif'' dari argumen diagonalisasi Cantor.

Kita dapat memberikan contoh yang lebih eksplisit tentang fungsi yang dapat
dihitung tetapi bukan rekursif primitif. Misalnya, biarkan notasi $g^n(x)$
menyatakan $g(g(\dots g(x)))$, dengan seluruhnya $n$ kemunculan~$g$; lalu
definisikan barisan fungsi $g_0,g_1,\dots$ dengan
\begin{eqnarray*}
g_0(x) & = & x+1 \\
g_{n + 1}(x) & = & g_n^x(x).
\end{eqnarray*}
Anda dapat memastikan bahwa setiap fungsi $g_n$ bersifat rekursif primitif.
Setiap fungsi berikutnya tumbuh jauh lebih cepat daripada fungsi sebelumnya;
$g_1(x)$ sama dengan $2x$, $g_2(x)$ sama dengan $2^x\cdot x$, dan $g_3(x)$
tumbuh kira-kira seperti menara eksponen yang terdiri atas $x$ buah~$2$.
Fungsi Ackermann--P\'eter pada dasarnya ialah fungsi $G(x)=g_x(x)$, dan dapat
ditunjukkan bahwa fungsi ini tumbuh lebih cepat daripada setiap fungsi
rekursif primitif.

Mari kembali ke masalah mengenumerasi fungsi-fungsi rekursif primitif. Ingat
bahwa kita telah menetapkan notasi simbolis bagi setiap fungsi rekursif
primitif; jadi, cukup mengenumerasi notasi-notasi tersebut. Kita dapat
menetapkan suatu bilangan asli $\#(F)$ bagi setiap notasi~$F$, secara rekursif,
sebagai berikut:
\begin{eqnarray*}
\#(0) & = & \langle 0 \rangle \\
\#(S) & = & \langle 1 \rangle \\
\#(\Proj{n}{i}) & = & \langle 2, n, i \rangle \\
\#(\fn{Comp}_{k,l}[H,G_0,\dots,G_{k-1}]) & = & \langle
3,k,l,\#(H),\#(G_0),\dots,\#(G_{k-1}) \rangle \\
\#(\fn{Rec}_l[G,H]) & = & \langle 4, l, \#(G), \#(H) \rangle.
\end{eqnarray*}
Di sini kita menggunakan fakta bahwa setiap barisan bilangan dapat dipandang
sebagai bilangan asli dengan menggunakan kode-kode dari bagian sebelumnya.
Hasilnya ialah bahwa setiap notasi diberi suatu bilangan asli. Tentu saja,
beberapa barisan (dan karenanya beberapa bilangan) tidak bersesuaian dengan
notasi; tetapi kita dapat menetapkan $f_i$ sebagai fungsi rekursif primitif
uner yang notasinya dikodekan oleh~$i$ jika $i$ mengodekan notasi semacam itu,
dan sebagai fungsi konstan~$0$ jika tidak. Hasil akhirnya ialah bahwa kita
mempunyai cara eksplisit untuk mengenumerasi fungsi-fungsi rekursif primitif
uner.

(Sebenarnya, beberapa fungsi, seperti fungsi nol konstan, akan muncul lebih
dari sekali dalam daftar. Ini bukan hanya akibat pengodean kita, tetapi juga
akibat fakta bahwa fungsi nol konstan mempunyai lebih dari satu notasi. Kita
akan melihat kemudian bahwa pengulangan ini tidak dapat dihindari secara
dapat dihitung; misalnya, tidak ada fungsi yang dapat dihitung untuk memutuskan
apakah suatu notasi tertentu merepresentasikan fungsi nol konstan.)

Sekarang kita dapat mengambil fungsi $g(x,y)$ yang diberikan oleh~$f_x(y)$,
dengan $f_x$ merujuk pada enumerasi yang baru saja dijelaskan. Bagaimana kita
tahu bahwa $g(x,y)$ dapat dihitung? Secara intuitif, hal ini jelas: untuk
menghitung $g(x,y)$, pertama-tama ``bongkar'' $x$ dan periksa apakah ia
merupakan notasi bagi fungsi uner. Jika demikian, hitung nilai fungsi itu pada
masukan~$y$.

\begin{tagblock}{TMs}
\begin{digress}
Anda mungkin sudah yakin bahwa (dengan sedikit usaha!) kita dapat menulis
program (misalnya, dalam Java atau C++) yang melakukan hal tersebut; lalu kita
dapat menggunakan tesis Church--Turing, yang menyatakan bahwa segala sesuatu
yang secara intuitif dapat dihitung dapat dihitung oleh mesin Turing.

Tentu saja, cara yang lebih langsung untuk menunjukkan bahwa $g(x,y)$ dapat
dihitung ialah menguraikan secara eksplisit suatu mesin Turing yang
menghitungnya. Secara khusus, ini akan menghindari tesis Church--Turing dan
daya tarik intuisi. Tidak lama lagi kita akan mempunyai cukup perangkat untuk
menunjukkan bahwa $g(x,y)$ dapat dihitung dengan menggunakan suatu model
komputasi yang dapat \emph{disimulasikan} pada mesin Turing, yakni fungsi
rekursif.
\end{digress}
\end{tagblock}

\end{document}
